% SGA 6, Exposición XII: fragmento español alineado con la fuente.
% Alcance: sección 2 completa, «Primeras reducciones».
% Fuente francesa: sga6_fr_r1/src/fr.tex, líneas 14219--14323.
% SHA-256 de la fuente francesa completa:
% 0D2158B9FE3B466309A92165727A0B191422E05FBAA0DACD5806E8322C54BA1A.

\section*{2. Primeras reducciones}

\textbf{2.1.} Apliquemos 1.3 a \(X/S\) y \(Y/S\): se obtienen un abierto no
vacío \(V\) de \(S\), un esquema íntegro \(S'\) y un morfismo finito y plano
\(S'\to V\), tales que toda componente irreducible \(X'_i\) de \(X_{S'}\) y
\(Y'_j\) de \(Y_{S'}\), dotada de la estructura reducida inducida, sea
íntegra sobre \(S'\). En virtud del descenso f.p.q.c. (EGA IV 2.7.1 (xiv) y
SGA 1 VIII 7.9), se puede sustituir \(S\) por \(S'\). Del siguiente diagrama
conmutativo:
\[
\begin{tikzcd}
X_{\red} \arrow[r,"f_{\red}"] \arrow[d,"i"'] & Y_{\red} \arrow[d,"j"]\\
X \arrow[r,"f"] & Y
\end{tikzcd}
\]
que implica el diagrama conmutativo
\[
\begin{tikzcd}
\Pic_{X_{\red}/S} & \Pic_{Y_{\red}/S} \arrow[l]\\
\Pic_{X/S} \arrow[u] & \Pic_{Y/S} \arrow[l] \arrow[u]
\end{tikzcd}
\]
y de sorites triviales sobre los morfismos de funtores representables por
morfismos cuasiafines (resultantes de EGA II 5.1.10 (ii) (v)), se deduce que
basta con demostrar las dos afirmaciones siguientes:
\begin{enumerate}[label=\Roman*.]
\item El teorema 1.1 cuando \(f\) es una inmersión (necesariamente
nilpotente).
\item El teorema 1.1 cuando \(X\) e \(Y\) son reducidos y toda componente
irreducible \(X_i\) de \(X\) y \(Y_i\) de \(Y\), dotada de la estructura
reducida inducida, es íntegra sobre \(S\).
\end{enumerate}

\textbf{2.2.} Estudiemos la afirmación II. Para todo \(i\), sea
\(j=\varphi(i)\) un índice tal que \(f(X_i)\subset Y_j\), y sea
\(\Sigma_j=\coprod_{\varphi(i)=j} X_i\). Se tienen entonces los dos
diagramas
\[
\begin{tikzcd}
Y & \coprod_j Y_j \arrow[l]\\
X \arrow[u] & \coprod_j \Sigma_j \arrow[l] \arrow[u]
\end{tikzcd}
\]
y
\[
% Canon SGA-SRC-001361: el diagrama español vigente ya lleva los objetivos Picard correctos.
\begin{tikzcd}
\Pic_{Y/S} \arrow[r,"\sgacircled{1}"] \arrow[d,"\sgacircled{2}"'] &
\Pic_{\coprod_jY_j/S}=\prod_j\Pic_{Y_j/S} \arrow[d,"\sgacircled{4}"]\\
\Pic_{X/S} \arrow[r,"\sgacircled{3}"] &
\Pic_{\coprod_j\Sigma_j/S}=\prod_j\Pic_{\Sigma_j/S}.
\end{tikzcd}
\]
Y, para establecer II, basta con ver que, tras restringir \(S\) si es
necesario, \(\sgacircled{1}\), \(\sgacircled{3}\) y
\(\sgacircled{4}\) son cuasiafines.

Consideremos primero \(\sgacircled{4}\). Sustituyendo \(Y\) por \(Y_j\) y
\(X\) por \(\Sigma_j\), nos reducimos al caso en que \(Y\) es íntegro sobre
\(S\). Como \(f\) es sobreyectivo, existe uno de los \(X_i\), sea \(X_0\),
que domina a \(Y\), y se tienen las dos sucesiones
\[
  X_0\hookrightarrow \coprod X_i\longrightarrow Y,
\]
y
\[
  \Pic_{X_0/S}\longleftarrow \Pic_{\coprod_i\sisi{X'_i}{X_i}/S}\longleftarrow \Pic_{Y/S}. % Canon SGA-SRC-001357
\]
Por tanto, si se sabe que los \(\underline{\Pic}_{X_i/S}\), y por consiguiente también
\(\Pic_{\coprod_iX_i/S}\), son representables, y que el morfismo compuesto
es cuasiafín, se concluye que \(\sgacircled{4}\) es cuasiafín. Teniendo en
cuenta la demostración de 1.2, para \(\sgacircled{4}\) nos reducimos, pues,
a demostrar II'. El teorema 1.1 cuando \(X\) e \(Y\) son íntegros sobre
\(S\).

Para \(\sgacircled{1}\) y \(\sgacircled{3}\), hay que tratar el caso de un
morfismo \(\coprod X_i\to X\), donde \(X\) es reducido, los \(X_i\) son las
componentes irreducibles de \(X\) y son íntegros sobre \(S\). Ahora bien, se
puede suponer que \(X\) tiene fibras conexas. En efecto, sea \(\eta\) el
punto geométrico de \(S\). Entonces \(X_\eta\) tiene como componentes
irreducibles los \((X_i)_\eta\) y, tras restringir \(S\) si es necesario, se
puede suponer que \(X_i\cap X_j=\varnothing\) si y solo si
\((X_i)_\eta\cap(X_j)_\eta=\varnothing\). Sean
\(Y_k=\bigcup_{i\in I(k)}X_i\) las componentes conexas de \(X\); entonces
\(f\) se descompone,
\(f=\coprod_k f_k:\coprod_k(\coprod_{i\in I(k)}X_i)\longrightarrow \sisi{Y_k}{\coprod_kY_k}\), y % Canon SGA-SRC-001358
de ello resulta la afirmación. Finalmente, una inducción inmediata nos
reduce al caso de un morfismo \(X_1\amalg X_2\to X\).
\(\sgacircled{1}\) y \(\sgacircled{3}\) resultan entonces, tras restringir
\(S\) si es necesario según EGA III 7.8.6, de la proposición siguiente.

\textbf{Proposición 2.3.} Sean \(S\) un esquema, \(g:Y\to S\) un morfismo
propio y \(Y_1,Y_2\) dos subesquemas cerrados definidos por los ideales
\(I_1,I_2\). Supongamos que
\begin{enumerate}[label=\alph*)]
\item \(Y,Y_1,Y_2,Y_1\cap Y_2\) son fielmente planos y de presentación
finita sobre \(S\).
\item \(g_*\mathcal O_Y=\mathcal O_S\) y
\((g_i)_*\mathcal O_{Y_i}=\mathcal O_S\) universalmente (donde
\(g_i=g|Y_i\), \(i=1,2\)).
\item \(I_1\cap I_2=0\).
\end{enumerate}
Entonces \(\Pic_{Y/S}\longrightarrow \Pic_{(Y_1\amalg Y_2)/S}\) es
representable por un morfismo afín de presentación finita.

En efecto, observemos primero que las hipótesis son estables bajo todo cambio
de base. Para (c), esto resulta de las dos sucesiones exactas:
\[
  0\longrightarrow I_1+I_2\longrightarrow\mathcal O_Y
  \longrightarrow\mathcal O_{Y_1\cap Y_2}\longrightarrow0
\]
y
\[
  0\longrightarrow I_1\cap I_2\longrightarrow I_1\oplus I_2
  \longrightarrow I_1+I_2\longrightarrow0~;
\]
la primera muestra que \(I_1+I_2\) es plano sobre \(S\); la segunda muestra
después que \(I_1\cap I_2\) es estable bajo cambio de base.

La afirmación es local para la topología f.p.p.f. Efectuando el cambio de
base \(Y_1\cap Y_2\to S\), se puede suponer que \(Y_1\cap Y_2/S\) posee una
sección \(\alpha\), que define secciones \(\alpha_i\) de \(Y_i\) y \(\beta\)
de \(Y\). Tras un cambio de base conveniente, nos reducimos a ver que todo
punto \(\mu:S\to\underline{\Pic}_{(Y_1\amalg Y_2)/S}\) tiene como imagen inversa un
subfuntor de \(\underline{\Pic}_{Y/S}\) que es representable por un esquema afín sobre
\(S\). Ahora bien, \(\mu\) está definido por un haz invertible \(L_i\) sobre
\(Y_i\), provisto de una rigidificación
\(u_i:\alpha_i^*L_i\xrightarrow{\sim}\mathcal O_S\), para \(i=1,2\). Del
mismo modo, un punto \(\lambda:S\to\underline{\Pic}_{Y/S}\) corresponde a un haz
invertible \(L\) sobre \(Y\), provisto de una rigidificación
\(u:\sisi{\alpha}{\beta}^*L\xrightarrow{\sim}\mathcal O_S\). La imagen de \(\lambda\) en % Canon SGA-SRC-001360
\(\Pic_{Y_1\amalg Y_2/S}\) es \(\mu\) si y solo si existe un isomorfismo de
haces rigidificados \(f_i^*(L)\simeq L_i\), donde
\(f_i:Y_i\hookrightarrow Y\) es la inclusión, para \(i=1,2\). Obsérvese que
estos isomorfismos definen un isomorfismo
\(\varphi:L_1|Y_1\cap Y_2\xrightarrow{\sim}L_2|Y_1\cap Y_2\) tal que
\(u_1=u_2\circ\varphi\). Recíprocamente, si \(\varphi\) es un isomorfismo
\(L_1|Y_1\cap Y_2\to L_2|Y_1\cap Y_2\) tal que
\(u_1=u_2\circ\varphi\), existe un haz invertible rigidificado \(L\) sobre
\(Y\), obtenido por pegado según \(\varphi\); explícitamente, \(L\) es el
núcleo del homomorfismo canónico
\((f_1)_*L_1\oplus(f_2)_*L_2\longrightarrow(f|Y_1\cap Y_2)_*L_2\) definido
por \(\varphi\). Este \(L\) corresponde entonces a un punto \(\lambda\) de
\(\Pic_{Y/S}\) que se envía a \(\mu\).

Finalmente, sea \(I\) el funtor definido, cuando \(T\) recorre los esquemas
sobre \(S\), por los conjuntos de isomorfismos
\(\varphi:(L_1)_T|(Y_1\cap Y_2)_T\xrightarrow{\sim}(L_2)_T|(Y_1\cap Y_2)_T\)
tales que \((u_1)_T=(u_2)_T\circ\varphi\). Acabamos de establecer un
isomorfismo entre la imagen inversa de \(\mu\) en \(\Pic_{Y/S}\) y el funtor
\(I\). Basta, pues, con aplicar el lema siguiente (donde se toma
\(X=Y_1\cap Y_2\));

\textbf{Lema 2.4.} Sea \(X\) un esquema plano, propio y de presentación
finita sobre \(S\). Sean \(L,M\) dos haces invertibles sobre \(X\), e
\(\Isom(L,M)\) el funtor de los isomorfismos \(L_T\to M_T\) (para \(T/S\)).
Entonces
\begin{enumerate}[label=\roman*)]
\item \(\underline{\Isom}(L,M)\) es representable por un esquema afín sobre \(S\), de
presentación finita.
\item Si, además, \(L\) y \(M\) están provistos de rigidificaciones
(relativamente a una sección dada de \(X\)), el subfuntor de \(\Isom(L,M)\)
formado por los isomorfismos compatibles con estas rigidificaciones es
representable por un cerrado de \(\Isom(L,M)\).
\end{enumerate}
En efecto, ii) está claro. En cuanto a i), recordemos que los funtores
\(T\mapsto\Hom(L_T,M_T)\) y \(T\mapsto\Hom(M_T,L_T)\) son representables por
fibrados vectoriales \(\mathbb V(E)\) y \(\mathbb V(F)\) (según EGA III
7.7.8). Entonces, si \(u\in\Hom(L_T,M_T)\) y
\(v\in\Hom(M_T,L_T)\), las condiciones \(u\circ v=1\) y \(v\circ u=1\)
definen un cerrado \(F\) de \(\mathbb V(E)\times_S\mathbb V(F)\), y está
claro que \(F\simeq \sisi{I}{\Isom(L,M)}\). % Canon SGA-SRC-001363

\textbf{2.5.} Finalmente, reduzcamos la afirmación II' a la afirmación
siguiente:

\(\mathrm{II}^*\). El teorema 1.1 cuando \(X\) e \(Y\) son íntegros sobre
\(S\) y la sucesión
\begin{equation}\tag{*}
  \mathcal O_Y
  \longrightarrow f_*\sisi{\mathcal O_Y}{\mathcal O_X} % Canon SGA-SRC-001364
  \rightrightarrows g_*\bigl(\mathcal O_{X\times_{\sisi{y}{Y}}X}\bigr)
\end{equation}
(donde \(g:X\times_{\sisi{y}{Y}}X\to Y\) es el morfismo estructural) es exacta.

De hecho, para esto basta con probar que existe una factorización de \(f\)
\[
  X\xrightarrow{f_0}Y_0\xrightarrow{f_1}Y_1\longrightarrow\cdots\xrightarrow{f_\eta}Y_\eta=Y
\]
tal que todo \(f_i\) satisfaga (*) y, tras restringir \(S\) si es necesario,
todo \(Y_i\) sea íntegro sobre \(S\).

Para ello, se pone \(Y_0=\Spec(f_*(\mathcal O_X))\), proviniendo \(f_0\) de
la factorización de Stein de \(f\) (EGA III 4.3.3), de modo que
\(f_{0*}(\mathcal O_X)\xleftarrow{\sim}\mathcal O_{Y_0}\), y que \(f_0\)
satisface efectivamente (*). Para factorizar del mismo modo \(Y_0\to Y\), se
utiliza el

\textbf{Lema 2.6.} Sean \(Y\) un esquema noetheriano,
\(f:X\to Y\) un epimorfismo finito (es decir, tal que
\(\mathcal O_Y\to f_*(\mathcal O_X)\) sea inyectivo); entonces \(f\) se
factoriza en una sucesión finita de morfismos finitos \(f_i\)
\(X\simeq Y_0\xrightarrow{f_1}Y_1\longrightarrow\cdots\longrightarrow Y_n=Y\)
que son epimorfismos efectivos (es decir, tales que la sucesión
\(\mathcal O_{Y_i}\longrightarrow (f_i)_*(\mathcal O_{Y_{i-1}})\rightrightarrows (g_i)_*\bigl(\mathcal O_{Y_{i-1}}\otimes_{\mathcal O_{Y_i}}\mathcal O_{Y_{i-1}}\bigr)\)
sea exacta).

Pongamos \(A=\mathcal O_Y\), \(B_0=f_*(\mathcal O_X)\), y definamos por
recurrencia
\(B_i=\Ker\bigl(B_{i-1}\rightrightarrows B_{i-1}\otimes_A B_{i-1}\bigr)\).
Se comprueba inmediatamente que
\(B_{i-1}\otimes_A B_{i-1}\xrightarrow{\sim}B_{i-1}\otimes_{\sisi{A_i}{B_i}}B_{i-1}\). % Canon SGA-SRC-001365
Queda por demostrar que, para \(i\) suficientemente grande, se tiene
\(B_i=A\). En efecto, consideremos la sucesión decreciente de
\(\sisi{B}{A}\)-módulos \(M_i=B_i/A\), y sea \(S_i=\Supp(M_i)\). La sucesión
decreciente de cerrados \(S_i\) de \(Y\) es estacionaria; sea \(S\) su valor
estacionario. Mostremos que \(S=\varnothing\). En caso contrario, sea \(s\)
un punto maximal de \(S\). Sustituyendo \(Y\) por
\(\Spec\mathcal O_{Y,s}\) si es necesario, se puede suponer que \(Y\) es
local, de punto cerrado \(s\). Por la elección de \(s\), \(M_i\) es
artiniano para \(i\) suficientemente grande, por lo que la sucesión
\((M_i)\) es estacionaria. Existe, pues, \(i\) tal que \(B_i=B_{i+1}\). Por
consiguiente, \(\Spec(\sisi{A_i}{B_i})\to Y\) es un monomorfismo y, por tanto, una
inmersión cerrada (Nakayama); como \(\sisi{B\subset A_i}{A\subset B_i}\), se tiene \(\sisi{B=A_i}{A=B_i}\), y
\(s\notin S\), lo cual es absurdo.
